\section{Model B: bounded total clearance}
%==========================================================================

\subsection{What is added}

Model A assumes alternative clearance grows without limit. Biologically it cannot:
carnitine is a finite cofactor and renal organic-acid transporters saturate. Model B lets
the second term saturate as well:

\begin{equation}
\dot m \;=\; s(t) \;-\; \theta\, g(m) \;-\; k_e\, h(m),
\qquad g(m) = \frac{m}{K_m + m}, \quad h(m) = \frac{m}{K_e + m}.
\label{eq:B}
\end{equation}

\subsection{Consequence: a load ceiling}

Total clearance $C(m) = \theta g(m) + k_e h(m)$ satisfies $C(0)=0$, is strictly increasing,
and is \emph{bounded}:
\begin{equation}
C(m) \;\longrightarrow\; \theta + k_e \qquad \text{as } m \to \infty .
\end{equation}

Equilibria are solutions of $C(m) = s$. Hence:

\begin{center}
\begin{tabular}{@{}ll@{}}
\toprule
$s < \theta + k_e$ & one equilibrium, stable \\
$s > \theta + k_e$ & \textbf{no equilibrium}; $\dot m > 0$ for all $m$, unbounded growth \\
\bottomrule
\end{tabular}
\end{center}

\begin{remark}[Terminology]
This is not a saddle-node bifurcation. No two equilibria collide; a single equilibrium
escapes to infinity as $s \to (\theta + k_e)^-$. A genuine saddle-node requires Model C.
\end{remark}

\begin{remark}[Clinical threshold precedes the mathematical one]
Clinically the crisis begins before the equilibrium ceases to exist, because the
equilibrium migrates to metabolite levels the patient does not tolerate. The operative
condition is $\spk > C(\bar m)$, with $\thc$ given by \eqref{eq:thcrit}.
\end{remark}

\subsection{Cost of the correction}

In Model A the two clearance terms have different functional forms, one saturating and one
linear, and are distinguished by that difference. In Model B both saturate, and they are
distinguished only by $K_m \neq K_e$, i.e.\ by saturating at different metabolite levels.

\textbf{The biologically more faithful model has weaker identifiability.} This trade-off
should be stated explicitly rather than discovered later.

\subsection{Status}

Model B is required for any statement about crisis onset and for the definition of $\thc$.
It is not required for the identifiability analysis of Section 3, and is therefore not the
working model for that purpose.

\vspace{0.5em}
\hrule
\vspace{1em}

%==========================================================================
